%=====================================================================
%  4. Experimental Design
%=====================================================================
\section{Experimental Design}
\label{sec:exp}

\subsection{Research Questions}

\begin{itemize}
\item[\textbf{RQ1 (adaptation).}] Does the multi-agent pipeline match
resources to learner profiles more accurately than single-pass generation?
\item[\textbf{RQ2 (coverage).}] What is the contribution of retrieval
grounding to core knowledge-point coverage?
\item[\textbf{RQ3 (reliability).}] How much does debate-based verification
reduce the hallucination rate, relative to no verification and to
self-refinement?
\item[\textbf{RQ4 (verification validity).}] Does the system's internal
hallucination estimate agree with independent human and external-model
judgments? \emph{This question exists because an internal judge is not, by
itself, sufficient evidence.}
\item[\textbf{RQ5 (perceived quality).}] How do human raters score the
generated bundles on correctness, fit, and usefulness?
\end{itemize}

\subsection{Materials}
\label{ssec:materials}

\emph{Topics.} We use \todo{14} topics drawn from an AI-for-industry
curriculum, ordered along a difficulty and application axis: supervised and
unsupervised learning foundations, deep learning and backpropagation, CNNs,
sequence modeling, Transformers and attention, LLMs and prompt engineering,
RAG, training and distributed tuning, model deployment and inference
acceleration, LLM evaluation and hallucination control, multi-agent system
development, manufacturing visual inspection, and predictive maintenance with
time-series analysis. Each topic carries a set of core knowledge points
$\mathcal{P}(t)$ used for coverage scoring, and a set of evidence queries used
for retrieval support measurement.

\emph{Learner profiles.} We construct \todo{12} profiles spanning vocational
diploma, BSc, MSc, and PhD students, zero-basis career changers, in-service
engineers, and corporate trainees. Each profile specifies educational
background, job target, per-subject theory scores, cognitive style, identified
weaknesses, learning goal, and preferred pace.

\emph{Cases.} Crossing profiles with topics under a stratified design yields
\todo{61} end-to-end cases. Each case records the profile, the topic, and the
difficulty label implied by the public rule of Section~\ref{ssec:diagnosis}.
A deterministic self-check confirms that the recorded labels agree with the
rule on \todo{100\%} of cases.

\emph{Knowledge base.} The domain corpus comprises \todo{4} source documents
segmented into \todo{56} retrieval chunks. Retrieval support is measured
separately from generation quality: for each evidence query we compute
bigram containment against all chunks and count support when containment
reaches $0.5$. Observed support is \todo{98.3\%} overall, indicating that low
coverage in generation cannot be attributed to missing evidence.

\subsection{Metrics}

\begin{description}
\item[Adaptation accuracy] the proportion of cases where the supervisor's
planned difficulty label equals the label implied by the public rule. Because
the rule is deterministic, this metric is externally checkable.
\item[Core knowledge-point coverage] $\frac{1}{|\mathcal{P}(t)|}\sum_{k \in
\mathcal{P}(t)} \mathbb{1}[k \text{ appears in the bundle}]$, measured by
keyword matching over the union of generated resources.
\item[Hallucination rate] the fraction of checked claims flagged as
unsupported, computed internally (debate agent), by external model judges, and
by human raters; we report all three and their agreement.
\item[Quality scores] reviewer ratings on accuracy, difficulty fit,
completeness, and cross-modal consistency (10-point scales).
\end{description}

\subsection{Baselines and Ablations}
\label{ssec:baselines}

To attribute effects to components rather than to the ensemble as a whole, we
compare against progressively weakened configurations
(Table~\ref{tab:baselines}). B1--B2 test whether orchestration is needed at
all; B3--B5 isolate the contribution of each verification mechanism.

\begin{table}[!t]
\centering
\caption{Evaluation configurations. All share the same LLM backend and
knowledge base; only the orchestration differs.}
\label{tab:baselines}
\begin{tabular}{@{}ll@{}}
\toprule
\textbf{ID} & \textbf{Configuration} \\
\midrule
B1 & Single LLM call, no profile, no retrieval \\
B2 & Single LLM call, profile in prompt, no retrieval \\
B3 & Full multi-agent, retrieval removed \\
B4 & Full multi-agent, debate stage removed \\
B5 & Full multi-agent, self-refinement replaces debate \\
B6 & Full multi-agent, review-gated retry disabled \\
\midrule
\textbf{Full} & \sysname{} with all components \\
\bottomrule
\end{tabular}
\end{table}

\subsection{Human Evaluation Protocol}
\label{ssec:human}

We recruit \todo{three} raters with \todo{graduate-level training in the
target domains}. Raters assess \todo{30} randomly sampled bundles, blind to
configuration, on three 5-point Likert items: factual correctness, difficulty
fit given the stated profile, and instructional usefulness. Each bundle is
rated by \todo{two} raters to permit agreement estimation. We report
percent agreement, Cohen's $\kappa$ \cite{cohen1960kappa}, and
Krippendorff's $\alpha$ \cite{krippendorff2011computing}. Raters additionally
mark any claim they judge factually wrong, which yields an independent
human hallucination count for RQ4.

\subsection{External Model Judgment}

To avoid the circularity inherent in having a system grade its own output with
the same underlying model, we re-judge all bundles with \todo{two} models from
a different family than the generating model, using a fixed rubric and no
access to the system's own verdicts. Agreement between internal, external, and
human judgments is reported as the primary validity evidence for RQ4.

\subsection{Statistical Analysis}

Because configurations are evaluated on the same cases, comparisons are
paired. We report mean differences with 95\% confidence intervals obtained by
\todo{10{,}000}-sample bootstrap, use the Wilcoxon signed-rank test when
normality is rejected by Shapiro--Wilk, and report Wilcoxon effect sizes
alongside Cohen's $d$ \cite{cohen2013statistical}. We apply Holm correction
for the family of comparisons in Table~\ref{tab:main}. All scripts are
released.

\subsection{Implementation Details}

All agents call \todo{the same commercial LLM endpoint} with per-agent
temperature settings in $[0.2,0.7]$; the verifier and reviewer use the lowest
temperatures. Orchestration is implemented as an explicit state graph over
HTTP; retrieval uses a persistent vector store with \todo{offline
deterministic embeddings}. Per-case wall-clock time and token cost are
reported in Table~\ref{tab:cost}.
